The autointerpretability process consists of five steps and yields both an interpretation and an autointerpretability score:

\begin{enumerate}
    \item On each of the first 50,000 lines of OpenWebText, take a 64-token sentence-fragment, and measure the feature's activation on each token of this fragment. Feature activations are rescaled to integer values between 0 and 10.
    \item Take the 20 fragments with the top activation scores and pass 5 of these to GPT-4, along with the rescaled per-token activations. Instruct GPT-4 to suggest an explanation for when the feature (or neuron) fires, resulting in an interpretation.
    \item Use GPT-3.5\footnote{While the process described in \citet{bills2023language} uses GPT-4 for the simulation step, we use GPT-3.5. This is because the simulation protocol requires the model's logprobs for scoring, and OpenAI's public API for GPT-3.5 (but not GPT-4) supports returning logprobs.} to simulate the feature across another 5 highly activating fragments and 5 randomly selected fragments (with non-zero variation) by asking it to provide the per-token activations.
    \item Compute the correlation of the simulated activations and the actual activations. This correlation is the autointerpretability score of the feature. The texts chosen for scoring a feature can be random text fragments, fragments chosen for containing a particularly high activation of that feature, or an even mixture of the two. We use a mixture of the two unless otherwise noted, also called `top-random' scoring.
    \item If, amongst the 50,000 fragments, there are fewer than 20 which contain non-zero variation in activation, then the feature is skipped entirely.
\end{enumerate}

Although the use of random fragments in Step 4 is ultimately preferable given a large enough sample size, the small sample sizes of a total of 640 tokens used for analysis mean that a random sample will likely not contain any highly activating examples for all but the most common features, making top-random scoring a desirable alternative.
